\noindent
\ProjectTitle{KōreroCare: Secure Mobile–Wearable Digital Phenotyping Infrastructure (2026 | Research)}

\noindent
\textit{Associated with Empathic Computing Laboratory}

\begin{itemize}
    \item Designed a multi-layer digital phenotyping architecture across Android (Kotlin), WearOS, and backend services to capture behavioral and physiological signals such as activity, heart rate, sleep, and contextual device usage data.

    \item Engineered high-uptime background sensing workflows using Foreground Services and WorkManager, enabling reliable passive data collection despite strict Android battery and background execution limits.

    \item Built resilient cross-device synchronization and buffering mechanisms to improve consistency between mobile and wearable streams and reduce data loss in real-world usage conditions.

    \item Developed secure and structured backend foundations using Go for service-layer design, with clear separation of ingestion, routing, and authentication responsibilities to support future microservice scaling.

    \item Integrated Keycloak-based authentication and identity management design to support secure token-based access control, role separation, and extensible researcher/user authentication workflows.

    \item Introduced Traefik as an API gateway and load-balancing layer to centralize routing, service exposure, and secure backend entry points in a containerized environment.

    \item Explored and incorporated Open Wearables as an integration layer for wearable ecosystem interoperability, helping normalize device data flows across heterogeneous wearable sources.

    \item Built validated ingestion APIs and data contracts for high-volume time-series streams, supporting structured handling of raw sensor data and downstream derived features.

    \item Defined schemas for both raw sensing data and processed feature outputs, creating a strong foundation for later machine learning, temporal analysis, and digital marker extraction.

    \item Solved multi-device timestamp alignment and synchronization challenges to improve the reliability of longitudinal data analysis across mobile and wearable sources.

    \item Designed observability and monitoring workflows to detect missing data, sensor failures, and real-world edge cases in production-like research deployments.

    \item Produced system architecture documentation and reproducibility-focused engineering artifacts to support collaboration, research continuity, and maintainable infrastructure evolution.
\end{itemize}

\noindent
\textit{Technologies:} Android, Kotlin, WearOS, Go, Keycloak, Traefik, Open Wearables, FastAPI, Pydantic, WorkManager, Foreground Services, Docker, Time-Series Data Pipelines, Digital Phenotyping, Observability
